% ============================================================
% CHAPTER 8
% ============================================================

\chapter{Entropy--Weighted Spectral Theory}
\label{chap:entropy-weighted-spectral}

\section{Motivation}

The spectral properties of lamphrodynamic operators encode geometric and
physical information about the entropy distribution $\SField$.
Eigenfunctions of entropy--weighted Laplacians describe resonant modes
of lamphrodynamic evolution and provide a bridge between field theory,
spectral geometry, and observable resonance phenomena (``cymatics'').

This chapter establishes the analytic foundations of entropy--weighted
spectral theory, proves self--adjointness of the main operators, and
derives spectral asymptotics under lamphrodynamic regularity
assumptions. These results will later inform Chapters~9 and~10, where we
relate spectral modes to cognitive and aesthetic structure.


\section{Entropy--Weighted Laplacian}

Let $(M,g)$ be a Riemannian manifold with entropy density $\SField$. The
entropy--weighted Laplacian is defined by
\[
  \Delta_{\SField} u := e^{-\SField}\nabla_a(e^{\SField}\nabla^a u).
\]

\begin{remark}
Entropy weighting biases diffusion so that entropy gradients induce
preferred propagation directions. This reflects lamphrodynamic flow of
semantic or physical information.
\end{remark}


\section{Self--Adjointness}

\begin{theorem}[Self--Adjointness]
\label{thm:self-adjoint}
Assume $\SField\in C^1(M)$ and $\SField$ is bounded from below. Then
$\Delta_{\SField}$ is essentially self--adjoint on $C_c^\infty(M)$ with
respect to the weighted measure $e^{\SField}\mathrm{dvol}_g$.
\end{theorem}

\begin{proof}[Sketch]
Integration by parts with weight $e^{\SField}$ shows symmetry. Standard
elliptic estimates and Friedrichs extension yield essential
self--adjointness.
\end{proof}


\section{Spectral Properties}

Consider the spectral problem
\[
  \Delta_{\SField} u = \lambda u.
\]

\begin{proposition}
Under the assumptions of \cref{thm:self-adjoint}, the spectrum of
$\Delta_{\SField}$ is real and bounded below, with discrete spectrum if
$M$ is compact.
\end{proposition}

\begin{proof}[Sketch]
Self--adjointness implies real spectrum. Compactness of $M$ yields
discreteness by Rellich compactness.
\end{proof}


\section{Heat Kernel Asymptotics}

The heat kernel associated to $\Delta_{\SField}$,
$K_{\SField}(t,x,y)$, satisfies
\[
  \partial_t K_{\SField} = \Delta_{\SField}K_{\SField}.
\]

\begin{theorem}
As $t\to 0^+$,
\[
  K_{\SField}(t,x,x) \sim (4\pi t)^{-n/2}\sum_{k=0}^\infty a_k(x)t^k,
\]
where $a_k$ depend on curvature of $g$ and derivatives of $\SField$.
\end{theorem}

\begin{proof}[Idea]
Adapt classical heat kernel asymptotics to weighted Laplacians using
perturbative expansions in $\SField$ and its derivatives.
\end{proof}


\section{Spectral Asymptotics}

\begin{proposition}
Let $\lambda_k$ denote eigenvalues of $\Delta_{\SField}$ on compact $M$.
Then
\[
  \lambda_k \sim C k^{2/n}
\]
for $k\to\infty$, with constant $C$ depending on $g$ and $\SField$.
\end{proposition}

\begin{proof}[Sketch]
Entropy weighting modifies volume density in Weyl’s law. Leading order
depends on $e^{\SField}\mathrm{dvol}_g$.
\end{proof}


\section{Physical Interpretation}

Eigenfunctions represent resonant lamphrodynamic patterns. High entropy
biases localization, while entropy gradients induce directional
propagation. Cymatic resonance arises when lamphrodynamic excitation
matches eigenvalues of $\Delta_{\SField}$.

\begin{remark}
Entropy weighting provides a physical mechanism for emergent resonance
patterns in lamphrodynamic systems, with potential applications to
cognition, neural oscillations, and aesthetic perception.
\end{remark}


\section{Open Problems}

\begin{enumerate}
\item Spectral gap estimates under entropy gradients.
\item Relation to lamphrodynamic phase transitions.
\item Experimental verification of entropy--weighted resonance.
\end{enumerate}

\newpage
